\section{All odd jet orders from two trees sharing a cyclic core}

\subsection*{Finding F-022: a uniform odd-order jet construction}
Let \(r\geq4\), \(R=r-1\), and \(\ell\geq1\). Put
\[
 b=2\ell+1,\qquad
 L_b(r)=2\binom{r+\ell-1}{\ell}-r.
                                                               \tag{F-022.1}
\]
The following statements hold in characteristic zero.
\begin{enumerate}
\item If \(n\geq L_b(r)\), every partition \(\lambda\vdash nd\)
with \(\ell(\lambda)\leq r\) and \(\lambda_2\leq b\)
satisfies \(K_n(\lambda)=a_\lambda(d[n])\).
\item If \(n\geq m\geq L_b(r)\), every such partition satisfies
\(q_\lambda(m,n)=K_n(\lambda)=a_\lambda(d[n])\), with
the internal padding \(X_{11}^{n-m}\operatorname{per}_m\).
\end{enumerate}
Both conclusions concern the entire multiplicity space, for
every outer degree \(d\). The coefficient maps have unit
Jacobian minors over \(\mathbb Z\), and their polynomial
pullbacks remain injective after arbitrary commutative base
change. This integral assertion is separate from the
characteristic-zero representation decomposition.
The cases \(\ell=1,2\) have sizes
\[
                         L_3(r)=r,\qquad L_5(r)=r^2.
                                                               \tag{F-022.2}
\]
The second recovers M-015 and its permanent variant F-021.
This is a checked generalization derived in this run, without
a priority claim.

\paragraph{The shared vertices and their labels.}
Write the variables as \(x,z_1,\ldots,z_R\), and let
\(\sigma=(1\ 2\ \cdots\ R)\). Take a central vertex
\(0\) and core vertices \(v_1,\ldots,v_R\).
For each nonnegative multi-index \(\alpha\) with
\(1\leq|\alpha|\leq\ell\), take an inward vertex
\(p_\alpha\) and an outward vertex \(q_\alpha\).
Make exactly the identifications
\[
                  p_{e_i}=v_i,\qquad
                  q_{e_j}=v_{\sigma^{-1}(j)}.
                                                               \tag{F-022.3}
\]
All vertices of degree at least two remain distinct, including
between the two families. The number of vertices is
\[
 1+2\left(\binom{R+\ell}{\ell}-1\right)-R=L_b(r).
\]
For \(|\alpha|\geq2\), remove one occurrence of its
largest supported index to obtain its inward parent:
\[
                 p_\alpha\longrightarrow p_{\alpha-e_j}
                    \quad\text{with label }z_j,
                    \qquad j=\max\operatorname{supp}\alpha.
\]
For \(|\beta|\geq2\), the outward edge instead removes
the smallest supported index in the reverse direction:
\[
                 q_{\beta-e_i}\longrightarrow q_\beta
                    \quad\text{with label }z_i,
                    \qquad i=\min\operatorname{supp}\beta.
\]
Finally put edges
\(v_i\longrightarrow0\) labeled \(z_i\) and
\(0\longrightarrow v_i\) labeled \(z_{\sigma(i)}\).
There are no other edges. Every nontrivial directed cycle
is a two-edge cycle through \(0\): inward paths strictly
decrease degree, outward paths strictly increase degree,
and neither family can return to the core after leaving it.

Let \(B\) be the positive-labeled adjacency matrix.
For the determinant use \(A_-=xI-B\), and for the
permanent use \(A_+=xI+B\). At \(x=1\), their values are
\[
 \det A_-=1-\sum_i z_i z_{\sigma(i)},\qquad
 \operatorname{per}A_+=1+\sum_i z_i z_{\sigma(i)}.
                                                               \tag{F-022.4}
\]
There is at most one nontrivial cycle in a cycle cover,
and every other vertex uses its diagonal loop. This proves
both identities directly.

\paragraph{Pure monomial cofactors.}
The cofactor of the central diagonal is one for both
polynomials, since deleting \(0\) leaves an acyclic graph.
For \(1\leq|\gamma|\leq\ell\), the path from \(0\)
to \(q_\gamma\) is unique and has product \(z^\gamma\).
It passes through the core vertex indexed by
\(\sigma^{-1}(\max\operatorname{supp}\gamma)\).

More generally, let \(\alpha,\beta\ne0\), with their
degrees at most \(\ell\), and suppose
\[
        \max\operatorname{supp}\beta
                      \ne\sigma(\min\operatorname{supp}\alpha).
                                                               \tag{F-022.5}
\]
The path from \(p_\alpha\) to \(q_\beta\) goes
inward to \(v_i\), through \(0\), and outward from
\(v_{\sigma^{-1}(j)}\), where
\(i=\min\operatorname{supp}\alpha\) and
\(j=\max\operatorname{supp}\beta\).
Condition (F-022.5) makes these two core vertices distinct.
The path is therefore simple and unique, and its product
is exactly \(z^{\alpha+\beta}\). It uses \(0\),
so no disjoint nontrivial cycle is available.

These paths give pure cofactors of both \(A_-\) and
\(A_+\). For the permanent this is the open-path minor
description in F-021, with no additional cycle terms because
the path occupies \(0\). For the determinant, differentiate
in the reverse entry to close the path into a cycle. If
the path has \(t\) edges, its negative labels contribute
\((-1)^t\), and the permutation sign of the resulting
\((t+1)\)-cycle is also \((-1)^t\). They cancel.
All unused loops are one and no other cycle is disjoint.
The resulting cofactor is the same positive monomial,
with coefficient one.

\begin{proof}[Coverage of every degree through \(2\ell\)]
Only degrees \(k\) with \(\ell<k\leq2\ell\) remain,
since the center-to-outward paths cover the smaller degrees.
Let \(\gamma\) have degree \(k\).

If its largest and smallest indices differ by at least two,
put a smallest index into \(\alpha\) and a largest into
\(\beta\), and distribute the remaining occurrences so
\(|\alpha|=\ell\), \(|\beta|=k-\ell\). Both degrees
are between one and \(\ell\). The smallest index \(i\)
is less than \(R\), and the largest \(j\) is at least
\(i+2\); hence \(j\ne\sigma(i)=i+1\).

If only one index occurs, split the multiplicity into
\(\ell\) and \(k-\ell\). Condition (F-022.5) follows
because \(\sigma\) has no fixed point.

The only remaining case uses adjacent indices \(t,t+1\)
with positive multiplicities \(a,b\). If \(a>\ell\),
put \(\ell\) low indices in \(\beta\), leaving
\(\alpha\) of degree \(k-\ell\). Its minimum and
\(\beta\)'s maximum are both \(t\), so the condition
holds. If \(b>\ell\), put \(\ell\) high indices in
\(\alpha\); its minimum and \(\beta\)'s maximum are
both \(t+1\), again giving the condition. These cases
cannot occur simultaneously because \(k\leq2\ell\).
Otherwise \(a,b\leq\ell\); put all high indices in
\(\alpha\) and all low indices in \(\beta\).
Now \(\sigma(t+1)\ne t\): if \(t+1<R\) its value
is \(t+2\), and at wraparound it is one whereas
\(t=R-1\geq2\). This uses \(R\geq3\), excluding a
two-cycle. Every case yields the required split and pure
cofactor. All degree bounds and the endpoint \(\ell=1\)
are included.
\end{proof}

\begin{proof}[Jet dominance and complete multiplicity spaces]
Vary a matrix entry by a linear form \(H\). The differential
of its determinant or permanent is \(H\) times the
corresponding cofactor. Taking \(H=x\) realizes all the
pure monomials just proved, and taking \(H=z_i\) supplies
every degree-\(2\ell+1\) monomial. Each target jet monomial
thus has an elementary coefficient direction whose derivative
is exactly that monomial. Distinct monomials require distinct
source coordinates. The chosen Jacobian minor is the identity
over \(\mathbb Z\).

The jet degree does not exceed the matrix size: the split
of \(z_i^{2\ell}\) gives a simple path of length
\(2\ell\), hence \(L_b(r)\geq2\ell+1=b\).
For a larger matrix append an \(xI\) block. Translation
by the appropriate coefficient vector in (F-022.4) gives,
on the chosen source subspace, \(F(t)=t+O(t^2)\)
integrally. The lowest nonzero homogeneous term of a
polynomial survives composition, proving injectivity after
arbitrary commutative base change and dominance over
\(\mathbb C\).

For the permanent, hold the diagonal at \(v_1\) fixed
as \(x\), and simultaneously permute its row and column
to the distinguished position \((1,1)\). This removes
none of the selected directions. The constant uses the
central diagonal; every other selected path is between
distinct vertices and uses an off-diagonal parameter.
This remains valid at \(\ell=1\), when no additional
tree vertices exist. Internal padding now gives exactly
\(x^{n-m}\operatorname{per}A_+\), whose dehomogenized
jets are unchanged.

All resulting pencil forms belong to the respective orbit
closures by endomorphism substitution. Every highest-weight
ideal vector with \(\lambda_2\leq b\) is a polynomial
in these jet coefficients by F-016. Pullback injectivity
forces it to be zero. Thus the whole ideal multiplicity
space vanishes, giving the two claimed ambient equalities.
The bounds also ensure that the \(r\)-variable space
fits in each matrix-entry ambient space, since \(L_b(r)\geq r\).
\end{proof}

\paragraph{An actual second-row gate at every even power.}
For every \(m\geq2\) and \(\ell\geq1\), an actual
internally padded permanent obstruction at
\(n\geq m^{2\ell}\) requires
\begin{equation}\label{eq:F022gate}
                              \lambda_2\geq2\ell+2.
\end{equation}
Indeed positive permanent multiplicity gives
\(r=\ell(\lambda)\leq m^2\). For \(r\leq3\),
O-006--O-007 already give full determinant multiplicity.
For \(r\geq4\), (F-022.2) proves
\(L_3(r)\leq m^2\) and \(L_5(r)\leq m^4\).
For \(\ell\geq3\), the proved M-013 estimate gives
\[
 L_{2\ell+1}(r)
    <2\binom{m^2+\ell-1}{\ell}
    \leq \frac58 m^{2\ell}\leq n.
\]
For completeness, its ratio to \(m^{2\ell}\) decreases
with \(m\geq2\), and its value at two is
\(2\binom{\ell+3}{3}/4^\ell\). This equals \(5/8\)
at \(\ell=3\), and decreases with ratio
\((\ell+4)/(4(\ell+1))\). Thus all parameter endpoints
are checked. The ambient determinant equality excludes
\(\lambda_2\leq2\ell+1\), proving the gate.
In particular \(n\geq m^2\) forces \(\lambda_2\geq4\),
and \(n\geq m^4\) forces \(\lambda_2\geq6\).

\paragraph{Grothendieck dictionary and limits.}
The coefficient targets are finite free quotients of the
polynomial algebra by powers of its augmentation ideal.
All matrices, cycle covers, cofactors, and chosen monomial
directions are integral. Their coordinate-algebra injections
are actual maps of affine schemes, universal under base
change. Passing to complete highest-weight spaces uses
F-016's Schur-functor support calculation and the comodule
formalism of SGA~3 I4.7.2--4.7.3 and SGA~6 VI3.3.
The extra information is the explicit integral injection,
not merely positivity in a representation \(K_0\).
Sharing core vertices allows cycles; the path/sign argument
proves precisely why they do not spoil the selected cofactors.
The results determine the stated full multiplicities and
exclude the stated actual gates, without any saturation
claim or comparison of still-undetermined evaluation ranks.
